\paragrafos{Μέτρηση γωνιών, σύγκριση και διχοτόμος}{Β.1.5, «Μέτρηση, σύγκριση και ισότητα γωνιών – Διχοτόμος γωνίας»}

\begin{anagke}
Στον κήπο υπάρχει μια γωνία φράχτη, και η Α3 θέλει να τη χωρίσει σε δύο ίσα παρτέρια με ένα μονοπάτι. «Πού βάζουμε το μονοπάτι;» ρωτάει ο Βασίλης. Η Αγγελική παίρνει το μοιρογνωμόνιο: «Πρώτα πρέπει να μετρήσουμε πόσο ανοίγει η γωνία». Βρίσκει $84$ μοίρες. «Το μισό είναι $42$», υπολογίζει η Πολυξένη. «Το μονοπάτι θα ξεκινήσει από την κορυφή και θα σχηματίζει $42$ μοίρες με κάθε πλευρά».

Ο Ηλίας, που φοράει ρολόι με δείκτες, παρατηρεί ότι οι δείκτες σχηματίζουν γωνία: «Στις $3{:}00$ οι δείκτες είναι κάθετοι, στις $6{:}00$ ευθεία». Η Ρήνη ρωτάει: «Μπορούμε να ξέρουμε πόσες μοίρες είναι η γωνία των δεικτών κάθε ώρα;». Ο Στέλιος δουλεύει με τον κύκλο: «Ολόκληρος ο κύκλος έχει $360$ μοίρες. Το ρολόι έχει $12$ ώρες. Κάθε ώρα είναι $360:12=30$ μοίρες».

Σε αυτή την παράγραφο μαθαίνουμε να μετράμε γωνίες με μοιρογνωμόνιο, να τις συγκρίνουμε και να τις χωρίζουμε σε δύο ίσα μέρη με τη διχοτόμο.
\end{anagke}

\begin{lexi}{μοίρα, μοιρογνωμόνιο, διχοτόμος}{μοίρα, μοιρογνωμόνιο, διχοτόμος}
Η \lexiword{μοίρα} είναι το «μέρος» (από το «μοίρα», μερίδιο): η μοίρα ($1^\circ$) είναι το $\frac{1}{360}$ του πλήρους κύκλου. Το \lexiword{μοιρογνωμόνιο} «γνωρίζει» (μετράει) τις μοίρες. Η \lexiword{διχοτόμος} «χωρίζει στα δύο» (δίχα + τέμνω) μια γωνία σε δύο ίσες γωνίες.
\end{lexi}

\begin{metafrasi}
Το $\frac{1}{360}$ του κύκλου. & $1^\circ$\\
Η υποδιαίρεση της μοίρας. & $1^\circ=60'$ και $1'=60''$\\
Η γωνία ανοίγει $84$ μοίρες. & $x\mathrm{Ο}y=84^\circ$\\
Η $\mathrm{Οδ}$ είναι διχοτόμος. & $x\mathrm{Ο}\delta=\delta\mathrm{Ο}y=\dfrac{x\mathrm{Ο}y}{2}$\\
Κάθε ώρα στο ρολόι. & $360^\circ:12=30^\circ$\\
\end{metafrasi}

\tmima{Η θεωρία}

\tmimasmall{Μέτρηση γωνιών}

Για να \textbf{μετρήσουμε} μια γωνία τη συγκρίνουμε με μια μονάδα, τη \textbf{μοίρα}. Μία μοίρα ($1^\circ$) είναι το $\frac{1}{360}$ του πλήρους κύκλου. Υποδιαιρείται σε \textbf{πρώτα λεπτά} και \textbf{δεύτερα λεπτά}: $1^\circ=60'$ και $1'=60''$. Για παράδειγμα $25^\circ30'=25{,}5^\circ$ και $48{,}25^\circ=48^\circ15'$, γιατί $0{,}25\cdot60=15$. Μετράμε με το \textbf{μοιρογνωμόνιο}: τοποθετούμε το κέντρο του στην κορυφή και τη μία πλευρά στο μηδέν, και διαβάζουμε πού περνά η άλλη πλευρά. Προσέχουμε σε ποια από τις δύο κλίμακες ξεκινά το μηδέν.

\tmimasmall{Σύγκριση και ισότητα}

Δύο γωνίες είναι \textbf{ίσες} όταν εφαρμόζουν η μία στην άλλη (έχουν το ίδιο μέτρο). Η μία είναι \textbf{μεγαλύτερη} από την άλλη όταν το μέτρο της είναι μεγαλύτερο, ή όταν η μικρή χωράει στο εσωτερικό της μεγάλης με κοινή κορυφή και μία κοινή πλευρά. Το μέτρο μιας γωνίας δεν εξαρτάται από το μήκος των πλευρών της (οι πλευρές είναι ημιευθείες).

\tmimasmall{Διχοτόμος}

\begin{orismos}{διχοτόμος γωνίας}
\textbf{Διχοτόμος} της γωνίας $x\mathrm{Ο}y$ είναι η ημιευθεία $\mathrm{Οδ}$ που βρίσκεται στο εσωτερικό της και τη χωρίζει σε δύο ίσες γωνίες: $x\mathrm{Ο}\delta=\delta\mathrm{Ο}y$. Τότε κάθε μία είναι το μισό της $x\mathrm{Ο}y$.
\end{orismos}

\begin{center}
\begin{tikzpicture}[scale=1.1]
\draw[thick] (3.5,0)--(0,0)--(2.6,2.2); \draw[gblue,thick] (0,0)--(3.2,1.1); \fill (0,0) circle (2pt);
\node[right] at (3.5,0) {$x$}; \node[above] at (2.6,2.2) {$y$}; \node[right,gblue] at (3.2,1.1) {$\delta$}; \node[below left] at (0,0) {$\mathrm{Ο}$};
\draw[gblue] (0.9,0) arc(0:19:0.9); \draw[gblue] (1.1,0.32) arc(19:40:1.1);
\draw (0.6,0.3) -- (0.6,0.3);
\end{tikzpicture}
\end{center}

Κάθε γωνία έχει μία μοναδική διχοτόμο. Την κατασκευάζουμε με μοιρογνωμόνιο (μετράμε, διαιρούμε με το $2$ και σημειώνουμε), ή με διαφανές χαρτί διπλώνοντας ώστε οι δύο πλευρές να συμπέσουν: η πτυχή είναι η διχοτόμος.

\begin{prosoxi}
\begin{itemize}[leftmargin=1.2em,itemsep=0.3ex,topsep=2pt]
\item Το μέτρο της γωνίας δεν εξαρτάται από το μήκος των πλευρών της: μια γωνία με μεγαλύτερο σχέδιο δεν είναι μεγαλύτερη.
\item Διαβάζω τη σωστή κλίμακα του μοιρογνωμονίου: αν η γωνία είναι οξεία το μέτρο της είναι μικρότερο από $90^\circ$.
\item Το $1^\circ=60'$ και όχι $100'$: $25^\circ30'=25{,}5^\circ$, όχι $25{,}30^\circ$.
\item Η διχοτόμος είναι ημιευθεία μέσα στη γωνία, όχι τμήμα που ενώνει δύο σημεία των πλευρών.
\end{itemize}
\end{prosoxi}

\tmima{Λυμένα παραδείγματα}

\paradeigma{Το μονοπάτι του κήπου}
\tixero Η γωνία του φράχτη είναι $84^\circ$.\par
\tithelo Η γωνία που σχηματίζει το μονοπάτι με κάθε πλευρά, αν ήθελαν να τη χωρίσουν σε δύο ίσα παρτέρια.\par
\possindeo Το μονοπάτι θα είναι η διχοτόμος. Κάθε γωνία είναι $84^\circ:2=42^\circ$.\par
\apantisi $42^\circ$.\par
\elegxos $42^\circ+42^\circ=84^\circ$.

\paradeigma{Οι δείκτες του ρολογιού}
\tixero Ο ωροδείκτης και ο λεπτοδείκτης στις $4{:}00$ και στις $5{:}00$.\par
\tithelo Η γωνία που σχηματίζουν.\par
\possindeo Κάθε ώρα αντιστοιχεί σε $360^\circ:12=30^\circ$. Στις $4{:}00$ ο λεπτοδείκτης δείχνει το $12$ και ο ωροδείκτης το $4$: $4\cdot30^\circ=120^\circ$. Στις $5{:}00$: $5\cdot30^\circ=150^\circ$.\par
\apantisi $120^\circ$ και $150^\circ$.\par
\elegxos Στις $6{:}00$ θα είναι $180^\circ$ (ευθεία).

\paradeigma{Μοίρες και λεπτά}
\tixero Οι γωνίες $25^\circ30'$ και $48{,}25^\circ$.\par
\tithelo Να τις γράψω στην άλλη μορφή.\par
\possindeo $30'=\frac{30}{60}^\circ=0{,}5^\circ$, άρα $25^\circ30'=25{,}5^\circ$. Και $0{,}25^\circ=0{,}25\cdot60'=15'$, άρα $48{,}25^\circ=48^\circ15'$.\par
\apantisi $25{,}5^\circ$ και $48^\circ15'$.\par
\elegxos $15'=\frac{15}{60}=0{,}25^\circ$.

\paradeigma{Σύγκριση γωνιών}
\tixero Οι γωνίες $42^\circ$, $4{,}5^\circ$, $38^\circ30'$, $38^\circ$.\par
\tithelo Η σειρά.\par
\possindeo Σε μοίρες: $42$, $4{,}5$, $38{,}5$, $38$.\par
\apantisi $4{,}5^\circ<38^\circ<38^\circ30'<42^\circ$.\par
\elegxos $30'=0{,}5^\circ$.

\paradeigma{Η διχοτόμος και το μέτρο}
\tixero Η $\mathrm{Οδ}$ είναι διχοτόμος της $x\mathrm{Ο}y$ και $x\mathrm{Ο}\delta=28^\circ$.\par
\tithelo Η $x\mathrm{Ο}y$.\par
\possindeo $\delta\mathrm{Ο}y=x\mathrm{Ο}\delta=28^\circ$, άρα $x\mathrm{Ο}y=28^\circ+28^\circ=56^\circ$.\par
\apantisi $56^\circ$.\par
\elegxos Η διχοτόμος της $56^\circ$ δίνει $28^\circ$.

\tmima{Ασκήσεις}

\zoni{Το κατάλαβα}{μοίρες, ρολόι, διχοτόμος}

\ex Σχεδίασε με μοιρογνωμόνιο γωνίες $40^\circ$, $75^\circ$ και $130^\circ$ και έλεγξε με το μέτρημα της διχοτόμου.

\ex Γράψε στην άλλη μορφή: α) $25^\circ30'$ σε μοίρες \quad β) $48{,}25^\circ$ σε μοίρες και λεπτά \quad γ) $3{,}5^\circ$ σε μοίρες και λεπτά \quad δ) $90^\circ$ σε λεπτά.

\ex Βάλε σε αύξουσα σειρά: $42^\circ$, $38^\circ30'$, $38^\circ$, $4{,}5^\circ$, $0{,}7^\circ$.

\ex Βρες τη διχοτόμο: α) $x\mathrm{Ο}y=70^\circ$ \quad β) $x\mathrm{Ο}y=124^\circ$ \quad γ) $x\mathrm{Ο}y=8^\circ30'$. Ποιο είναι το μέτρο κάθε μισής γωνίας;

\ex Η $\mathrm{Οδ}$ είναι διχοτόμος της $x\mathrm{Ο}y$ και $x\mathrm{Ο}\delta=28^\circ$. Βρες την $x\mathrm{Ο}y$.

\ex Βρες τη γωνία των δεικτών του ρολογιού στις: α) $3{:}00$ \quad β) $6{:}00$ \quad γ) $9{:}00$ \quad δ) $2{:}00$ \quad ε) $5{:}00$ \quad στ) $12{:}00$. (Πάρε τη μικρότερη γωνία.)

\ex Σωστό ή λάθος; α) $1^\circ=60'$. \quad β) Η διχοτόμος χωρίζει τη γωνία σε δύο άνισα μέρη. \quad γ) Δύο γωνίες είναι ίσες αν εφαρμόζουν. \quad δ) Η διχοτόμος είναι ημιευθεία. \quad ε) $1'=60''$.

\zoni{Το χρησιμοποιώ}{κήπος, ρολόι, πίτσα}

\ex \textbf{Το μονοπάτι.} Η γωνία του φράχτη είναι $84^\circ$. Με ποια γωνία ξεκινά το μονοπάτι που τη χωρίζει σε δύο ίσα παρτέρια;

\ex \textbf{Το ρολόι.} Πότε (ποια ώρα, με τον λεπτοδείκτη στο $12$) οι δείκτες σχηματίζουν γωνία $150^\circ$; Υπάρχουν δύο ώρες;

\ex \textbf{Η πίτσα.} Μια πίτσα κόβεται σε $8$ ίσες φέτες. Πόση γωνία έχει η κορυφή κάθε φέτας; Αν κόβεται σε $6$ φέτες;

\ex \textbf{Ο ανεμόμυλος.} Ένας ανεμόμυλος έχει $3$ πτερύγια με ίσες γωνίες μεταξύ τους. Πόση γωνία σχηματίζουν δύο διαδοχικά;

\ex \textbf{Το μοιρογνωμόνιο.} Ο Στέλιος διάβασε $140^\circ$ σε μια γωνία που φαίνεται οξεία. Τι πιθανότατα έκανε; Ποιο είναι το σωστό;

\zoni{Το σκέφτομαι}{γιατί $360$;}

\ex Οι Βαβυλώνιοι χωρίζουν τον κύκλο σε $360$ μοίρες. Γιατί ο $360$ είναι βολικός αριθμός (σκέψου τους διαιρέτες του);

\ex Αν διπλασιάσεις το μήκος των πλευρών μιας γωνίας, αλλάζει το μέτρο της; Εξήγησε.

\ex Γιατί κάθε γωνία έχει μία και μοναδική διχοτόμο;

\ex Η διχοτόμος μιας ορθής γωνίας ($90^\circ$) πόσες μοίρες σχηματίζει με κάθε πλευρά;

\zoni{Για το σπίτι}{δίπλωσε το χαρτί}

\ex Κόψε μια γωνία από χαρτί και βρες τη διχοτόμο διπλώνοντάς την. Μέτρησε με μοιρογνωμόνιο.

\ex Μέτρησε με μοιρογνωμόνιο τρεις γωνίες στο σπίτι και γράψε τα μέτρα τους.
